%------------------------------------------------------------------------------%
\section{Teste da Integral}

Nesta seção e nas próximas vamos apresentar alguns testes que podem ser utilizados
para verificar se séries com termos não-negativos são convergentes. Começamos
definindo esse caso particular de séries.

%------------------------------------------------------------------------------%
\begin{definicao}{Série com Termos Não Negativos}{def:serie-termos-nao-negativos}
  Uma série $\sum a_n$ onde todos os termos da sequência $a_n$ são não negativos,
  isso é,
  \[
    a_n \geq 0 \qquad \forall n
  \]
  é chamada de \emph{Série de Termos Não Negativos}\index{Série!termos não negativos}.
\end{definicao}
%------------------------------------------------------------------------------%

Note que se quisermos estudar uma série com ternos não positivos \(a_n\leq 0\)
podemos remover todos os sinais negativos e estudar a série de termos
não negativos correspondente. Uma das razões para começarmos o estudo
da convergência das séries por esse caso particular é que podemos usar
o Teorema da Sequência Monotônica~\ref{teo:sequencia-monotona} na
sequência das somas parciais o que produz o seguinte corolário.

%------------------------------------------------------------------------------%
\begin{corolario}{Convergência de Série com Termos Não Negativos}
                 {cor:convergencia-serie-termos-nao-negativos}
    Uma série $\sum a_n$ de Termos Não Negativos, $a_n \geq 0$,
    é \emph{convergente} se, e somente se, sua sequência de somas
    parciais é uma sequência \emph{limitada}.
\end{corolario}
%------------------------------------------------------------------------------%
\begin{proof}
  Basta ver que a sequência de somas parciais é crescente,
  já que a sequência do termo geral é não negativa.
  Assim, $S_n$ é monotônica e portanto convergente.
  Por outro lado, se a série é convergente,
  então a sequência $S_n$ é convergente e, portanto, limitada.
\end{proof}
%------------------------------------------------------------------------------%

Os próximos exemplos ilustram o uso desse resultado.
No início do Capítulo definimos a Série do Inverso do Quadrado
(Definição~\ref{def:serie-inverso-quadrado}) e apresentamos
sua soma sem demonstração, vamos agora mostrar que ela converge.

%------------------------------------------------------------------------------%
\begin{example}[ex:convergencia-serie-inverso-quadrado]
  Mostre que a
  \emph{Série do Inverso do Quadrado}\index{Série!inverso do quadrado}
  é convergente.
  \solution
  Verificamos facilmente que todos os termos são positivos e
  construímos a sequência de somas parciais
  \[
    S_n = \sum_{k=1}^n \frac{1}{k^2}
        = 1 + \frac{1}{4}
            + \frac{1}{9}
            + \frac{1}{16}
            + \frac{1}{25}
            + \cdots
            + \frac{1}{n^2}    \,.
  \]
  Começamos manipulando $S_n$
  \begin{align*}
    S_n & = 1 + \frac{1}{4}
        + \frac{1}{9}
        + \frac{1}{16}
        + \frac{1}{25}
        + \cdots
        + \frac{1}{n^2}    \\[3mm]
        & = 1 + \frac{1}{2\cdot 2}
        + \frac{1}{3\cdot 3}
        + \frac{1}{4\cdot 4}
        + \frac{1}{5\cdot 5}
        + \cdots
        + \frac{1}{nn}     \\[3mm]
        & < 1 + \frac{1}{1\cdot 2}
        + \frac{1}{2\cdot 3}
        + \frac{1}{3\cdot 4}
        + \frac{1}{4\cdot 5}
        + \cdots
        + \frac{1}{(n-1)n} \\[3mm]
        & = 1 + \sigma_n
  \end{align*}
  onde $\sigma_n$ são todos os termos exceto o primeiro
  \begin{align*}
    \sigma_n & = \frac{1}{1\cdot 2}
               + \frac{1}{2\cdot 3}
               + \frac{1}{3\cdot 4}
               + \frac{1}{4\cdot 5}
               + \cdots
               + \frac{1}{(n-1)n} \\[3mm]
             & = \sum_{k=2}^n \frac{1}{(k-1)k}            \,,
  \end{align*}
  ou seja, $\sigma_n$ é a sequência de somas parciais da série com termos
  \[
    a_k = \frac{1}{(k-1)k}         \,.
  \]
  Decompondo $a_k$ em frações parciais
  \[
    a_k = \frac{1}{\,k-1\,} - \frac{1}{\,k\,}
  \]
  verificamos que essa é uma Série Telescópica e portanto podemos calcular seu limite
  \begin{align*}
    \sigma & = \lim_{n\to\infty} \sigma_n  \\[3mm]
           & = \lim_{n\to\infty} \sum_{k=2}^n \left( \frac{1}{\,k-1\,} - \frac{1}{\,k\,}\right) \\[3mm]
           & = \lim_{n\to\infty} \left(\frac{1}{\,2-1\,} - \frac{1}{\,n\,}\right) \\[3mm]
           & = \lim_{n\to\infty} 1 - \frac{1}{\,n\,} \\[3mm]
           & = 1        \,.
  \end{align*}
  Retornando a $S_n$, para todo $n\geq 2$ podemos escrever
  \[
    S_n < 1 + \sigma_n
  \]
  calculando o limite nos dois lados da desigualdade verificamos que
  \[
    \sum_{n=1}^{\infty} \frac{1}{n^2} \;<\; 2          \,,
  \]
  ou seja, a série é limitada.
  Portanto pelo Corolário~\ref{cor:convergencia-serie-termos-nao-negativos}
  ela é convergente.
\end{example}
%------------------------------------------------------------------------------%

Vamos também mostrar que a Série do Inverso do Fatorial~\ref{def:serie-inverso-fatorial}
é convergente.

%------------------------------------------------------------------------------%
\begin{example}[ex:convergencia-serie-inverso-fatorial]
  Mostre que a
  \emph{Série do Inverso do Fatorial}\index{Série!inverso do fatorial}
  é convergente.
  \solution
  Também nesse caso é fácil verificar que os termos da série são todos
  positivos e começamos manipulando suas somas parciais
  \begin{align*}
    S_n & = \sum_{k=0}^n \frac{1}{k!} \\[3mm]
        & = 1 + 1
        + \frac{1}{2}
        + \frac{1}{3\cdot2}
        + \frac{1}{4\cdot 3\cdot 2}
        + \frac{1}{5\cdot 4\cdot 3\cdot 2}
        + \cdots + \frac{1}{n\cdot (n-1)\cdots 3\cdot 2} \\[3mm]
        & < 1 + 1
        + \frac{1}{2}
        + \frac{1}{2\cdot2}
        + \frac{1}{2\cdot 2\cdot 2}
        + \frac{1}{2\cdot 2\cdot 2\cdot 2}
        + \cdots
        + \frac{1}{2\cdot 2\cdots 2\cdot 2} \\[3mm]
        & = 1 + 1
        + \frac{1}{2}
        + \frac{1}{2^2}
        + \frac{1}{2^3}
        + \frac{1}{2^4}
        + \cdots
        + \frac{1}{2^{n-1}}          \,.
  \end{align*}
  Observe que, removendo o primeiro termo, a última expressão é a
  $n$-ésima soma parcial da série geométrica
  com $\alpha=1$ e $r=\sfrac{1}{2}$,
  \[
    \sigma_n =  1
             + \frac{1}{2}
             + \frac{1}{2^2}
             + \frac{1}{2^3}
             + \frac{1}{2^4}
             + \cdots
             + \frac{1}{2^{n-1}}          \,,
  \]
  cuja soma é
  \[
    \sigma = \frac{\alpha}{1-r}
           = \frac{1}{1-\sfrac{1}{2}}
           = \frac{1}{\sfrac{1}{2}}
           = 2            \,.
  \]
  Retornando a $S_n$, temos que, para todo $n\geq 2$ vale a seguinte desigualdade
  \[
    S_n < 1 + \sigma_n      \,.
  \]
  Calculando o limite dos dois lados da desigualdade verificamos que
  \[
    \sum_{n=0}^\infty \frac{1}{\,n!\,} < 3           \,.
  \]
  Como a série só tem termos não negativos e é limitada, usamos o
  Corolário~\ref{cor:convergencia-serie-termos-nao-negativos}
  para concluir que a série é convergente.
\end{example}
%------------------------------------------------------------------------------%

%------------------------------------------------------------------------------%
\begin{figure}
  \centering
  \begin{subfigure}{\linewidth}
    \centering
    \vspace{2mm}
    % Title: gl2ps_renderer figure
% Creator: GL2PS 1.4.2, (C) 1999-2020 C. Geuzaine
% For: Octave
% CreationDate: Fri Sep 17 21:56:23 2021
\setlength{\unitlength}{1pt}
\begin{picture}(0,0)
\includegraphics[width=16cm,height=5cm]{fig_teste_integral_right}
\end{picture}%
\begin{picture}(453,141)(0,0)
\fontsize{10}{0}\selectfont\put(120.945,10.9531){\makebox(0,0)[t]{\textcolor[rgb]{0.15,0.15,0.15}{{$N\!-\!1$}}}}
\fontsize{10}{0}\selectfont\put(156.22,10.9531){\makebox(0,0)[t]{\textcolor[rgb]{0.15,0.15,0.15}{{$N$}}}}
\fontsize{10}{0}\selectfont\put(191.496,10.9531){\makebox(0,0)[t]{\textcolor[rgb]{0.15,0.15,0.15}{{$N\!+\!1$}}}}
\fontsize{10}{0}\selectfont\put(262.047,10.9531){\makebox(0,0)[t]{\textcolor[rgb]{0.15,0.15,0.15}{{$N\!+\!3$}}}}
\fontsize{10}{0}\selectfont\put(332.598,10.9531){\makebox(0,0)[t]{\textcolor[rgb]{0.15,0.15,0.15}{{$N\!+\!5$}}}}
\fontsize{10}{0}\selectfont\put(403.15,10.9531){\makebox(0,0)[t]{\textcolor[rgb]{0.15,0.15,0.15}{{$N\!+\!7$}}}}
\end{picture}

    \caption{Construindo os retângulos a direita a série é maior do que a integral.}
    \label{fig:teste_integral_right}
  \end{subfigure}\\[2mm]
  \begin{subfigure}{\linewidth}
    \centering
    \input{fig_teste_integral_left}
    \caption{Construindo os retângulos a esquerda a série é menor do que a integral.}
    \label{fig:teste_integral_left}
  \end{subfigure}\\
  \caption{Ilustração do teste da integral.}
  \label{fig:teste_integral}
\end{figure}
%------------------------------------------------------------------------------%

Vamos agora apresentar o teste da integral que compara a soma da série com a Integral
Imprópria~\ref{def:integral-impropria-1} de uma função~$f$.
Porém, antes de apresentar esse resultado, vamos explorar as ideias envolvidas
analisando a Figura~\ref{fig:teste_integral}.
Estamos assumindo que, pelo menos a partir de algum ponto $n\geq N$,
os termos da série sejam dados pela função, $a_n=f(n)$,
e que a função é contínua e decrescente para  $x\geq N$.
Vamos representar os termos da série por retângulos de largura um e altura $a_n$,
que podem ser construídos a direita ou a esquerda do ponto $x=n$.
Como $f$ é decrescente, se os retângulos estão na direita
como no Gráfico~\ref{fig:teste_integral_right},
a área do retângulo é maior do que a integral de $f$ em cada intervalo $[n, n+1]$.
Enquanto que, se desenharmos o retângulo à esquerda
como no Gráfico~\ref{fig:teste_integral_left},
a integral será maior do que a área do retângulo.
O próximo teorema usa essas relações para associar a convergência da integral
imprópria com a existência da soma da série.

%------------------------------------------------------------------------------%
\begin{teorema}{Teste da Integral}{teo:teste-integral}
  Seja\index{Teste!integral}
  $a_n$ uma sequência tal que para todo $n\geq N$, onde $N$ é um número natural,
  $a_n\geq 0$. Seja $f$ uma função contínua, positiva e decrescente para $x\geq N$,
  satisfazendo $f(n)=a_n$, para todo $n\geq N$. Então, a série
  \[
    \sum_{n=N}^\infty a_n
  \]
  converge se, e somente se, a integral
  \[
    \int_N^\infty f(x)\,dx
  \]
  converge.
\end{teorema}
%------------------------------------------------------------------------------%
\begin{proof}
\emph{Prova da Ida}

  Vamos mostrar primeiro que a convergência da série garante a convergência
  da integral.
  Seja $n>N$ e considere a partição do intervalo $[N, n]$ dada por
  \[
    [N, N+1] \cup [N+1, N+2]
             \cup \cdots
             \cup [n-1, n]    \,.
  \]
  Como cada intervalo tem comprimento igual a $1$,
  a função $f$ é decrescente e estamos considerando os retângulos a direita do
  ponto $x=n$, segue que a soma parcial
  \[
    S_{N,n-1} = a_N + a_{N+1} + \cdots + a_{n-1}
  \]
  nos dá um limite superior para área abaixo do gráfico da função $f$
  dentro do intervalo $[N, n]$, isso é,
  \[
    \int_N^n f(x)\,dx \leq S_{N,n-1}  \,.
  \]
  Como a série converge a sequência de somas parciais $S_{N,n-1}$
  converge para algum valor $L$.
  Além disso, como $a_n$ é não negativo a sequência $S_{N,n-1}$ é crescente,
  então para todo $n > N$
  \[
    I_n = \int_N^n f(x)\,dx \leq  L   \,,
  \]
  ou seja, a sequência $I_n$ é uma sequência limitada,
  além disso como $f$ é positiva $I_n$ é crescente e podemos
  concluir que a sequência $I_n$ é convergente e portanto
  a integral imprópria
  \[
    \int_N^\infty f(x)\,dx
  \]
  é convergente.

%------------------------------------------------------------------------------%
\emph{Verificando a Volta}

  Agora mostraremos que se a integral converge a série também converge.
  Para isso vamos construir os retângulos que representam os termos da
  série à esquerda do ponto $x=n$, nesse caso a soma parcial correspondente
  ao intervalo $[N, n]$ é dada por
  \[
    S_{N+1,n} = a_{N+1} + a_{N+2} + \cdots + a_{n}      \,.
  \]
  Usando novamente o fato que a função é decrescente,
  essa soma parcial nos dá um limite inferior para a integral, ou seja
  \[
    S_{N+1,n} \leq \int_N^n f(x)\,dx           \,.
  \]
  Somando $a_N$ dos dois lados temos que
  \[
    S_{N,n} = a_N + S_{N+1,n} \leq a_N + \int_N^n f(x)\, dx       \,.
  \]
  Como a função $f$ é positiva e a integral imprópria é convergente
  \[
    \int_N^\infty f(x)\,dx = L
  \]
  temos, para todo $n>N$,  que
  \[
    S_{N,n} \leq L     \,.
  \]
  Como $a_k$ é não negativo e a série é limitada
  usando o Corolário~\ref{cor:convergencia-serie-termos-nao-negativos}
  concluímos que a série é convergente.
\end{proof}
%------------------------------------------------------------------------------%

Note que também podemos reescrever a tese do teorema acima como:
a série
\[
  \sum_{n=N}^\infty a_n
\]
diverge se, e somente se, a integral
\[
  \int_N^\infty f(x)\,dx
\]
diverge. O próximo exemplo aplica essa versão do teorema.

%------------------------------------------------------------------------------%
\begin{example}
  Verificar a convergência da integral
  \(
    \int_{1}^\infty \frac{\,dx\,}{x}
  \).
  \solution
  Verificamos que a Série Harmônica
  \[
    \sum_{n=1}^\infty \frac{1}{\,n\,}
  \]
  diverge, assim, pela observação acima, podemos concluir que a integral
  \[
    \int_{1}^\infty \frac{\,dx\,}{x}
  \]
  também diverge.
\end{example}
%------------------------------------------------------------------------------%

É importante destacar que o valor a integral \emph{não é uma boa aproximação}
para a soma da série.
Podemos observar isso na Figura~\ref{fig:teste_integral_comparacao}
comparando a área sob o gráfico da função $f(x)=x^{-2}$
entre os pontos $n$ e $n+1$ que é data pela integral
\[
  I_n = \int_n^{n+1} \frac{\,dx\,}{x^2} = \frac{1}{\,n^2+n\,}
\]
com os termos da série
\[
  a_n = \frac{1}{\,n^2\,}
\]
representados pelos retângulos.

%------------------------------------------------------------------------------%
\begin{figure}
  \centering
  % Title: gl2ps_renderer figure
% Creator: GL2PS 1.4.2, (C) 1999-2020 C. Geuzaine
% For: Octave
% CreationDate: Sat Sep 18 17:28:01 2021
\setlength{\unitlength}{1pt}
\begin{picture}(0,0)
\includegraphics[width=14cm,height=5cm]{fig_teste_integral_comparacao}
\end{picture}%
\begin{picture}(396,141)(0,0)
\fontsize{10}{0}\selectfont\put(56.0551,9.94382){\makebox(0,0)[t]{\textcolor[rgb]{0.15,0.15,0.15}{{1}}}}
\fontsize{10}{0}\selectfont\put(125.504,9.94382){\makebox(0,0)[t]{\textcolor[rgb]{0.15,0.15,0.15}{{3}}}}
\fontsize{10}{0}\selectfont\put(194.953,9.94382){\makebox(0,0)[t]{\textcolor[rgb]{0.15,0.15,0.15}{{5}}}}
\fontsize{10}{0}\selectfont\put(264.402,9.94382){\makebox(0,0)[t]{\textcolor[rgb]{0.15,0.15,0.15}{{7}}}}
\fontsize{10}{0}\selectfont\put(333.85,9.94382){\makebox(0,0)[t]{\textcolor[rgb]{0.15,0.15,0.15}{{9}}}}
\fontsize{10}{0}\selectfont\put(16.3318,72.9685){\makebox(0,0)[r]{\textcolor[rgb]{0.15,0.15,0.15}{{0,5}}}}
\fontsize{10}{0}\selectfont\put(16.3318,128.48){\makebox(0,0)[r]{\textcolor[rgb]{0.15,0.15,0.15}{{1,0}}}}
\end{picture}

  \caption{Comparação do valor da integral com a soma da série,
    mostrando que a integral não é uma boa aproximação para a soma.}
  \label{fig:teste_integral_comparacao}
\end{figure}
%------------------------------------------------------------------------------%

A tabela a seguir reproduz a mesma comparação apresentando os valores numéricos
dos sete primeiros termos e sua soma.
Note que os termos são significativamente diferentes no início e só
começam a se aproximar posteriormente. Note também que o erro acumulado
no início nunca é compensado não importando o número de termos somado.

\begin{ntabular}{cccc}[0.6]
  $n$   & $a_n$   & $I_n$   & $a_n-I_n$  \\ \hline \hline
  1     & 1,000   & 0,500   & 0,500      \\ \hline
  2     & 0,250   & 0,167   & 0,083      \\ \hline
  3     & 0,111   & 0,083   & 0,028      \\ \hline
  4     & 0,062   & 0,050   & 0,012      \\ \hline
  5     & 0,040   & 0,033   & 0,007      \\ \hline
  6     & 0,028   & 0,024   & 0,004      \\ \hline
  7     & 0,020   & 0,018   & 0,003      \\ \hline \hline
  Soma  & 1,512   & 0,875   & 0,637
\end{ntabular}

Mais a frente vamos apresentar o Teorema~\ref{teo:estimativa-erro},
que nos mostra como podemos usar o valor da integral para obter estimativas
de erro para a aproximação da soma da série por suas somas parciais.

Podemos agora utilizar o Teste da Integral para verificar a convergência
de uma classe importante de séries, que chamamos de Serie~$p$. Primeiro
apresentamos a definição dessas séries e depois aplicamos o teste para
verificar em quais condições elas convergem.

%------------------------------------------------------------------------------%
\begin{definicao}{Série $p$}{def:serie-p}
  Uma série com a forma
  \(\;
    \sum_{n=1}^\infty \frac{1}{\,n^p\,}
  \;\)
  é chamada de \emph{Série $p$}\index{Série!$p$} ou \emph{$p$-série}.
\end{definicao}
%------------------------------------------------------------------------------%

Essas séries são importantes pois existe uma forma simples para verificar
quando elas são convergentes. O próximo resultado apresenta esse critério.

%------------------------------------------------------------------------------%
\begin{proposicao}{Convergência da Série $p$}{pro:convergencia-serie-p}
  Uma série $p$ converge se $p>1$ e diverge se $p\leq 1$.
\end{proposicao}
%------------------------------------------------------------------------------%
%------------------------------------------------------------------------------%
\begin{proof}
  Quando $p=1$ temos a série harmônica que diverge, como já vimos.

  Suponha então que $p\neq 1$, nesse caso temos
  \begin{align}
    \int_1^\infty\frac{\,dx\,}{x^p}
         & = \lim_{b\to \infty}\left[\;\frac{x^{-p+1}}{-p+1}\evalat{1}{b} \; \right] \nonumber \\[3mm]
         & = \frac{1}{1-p} \;\lim_{b\to \infty}\left(\frac{1}{b^{p-1}}-1\right)   \,.
         \label{eq:int1divxp}
  \end{align}
  Quando $p>1$ temos que $p-1>0$, então
  \[
    \lim_{b\to\infty} \frac{1}{\,b^{p-1}\,} = 0      \,.
  \]
  Assim, da equação \eqref{eq:int1divxp} temos que a integral converge e
  \[
    \int_1^\infty\frac{\,dx\,}{x^p}=\frac{1}{\,p-1\,}     \,.
  \]
  Concluímos então que para $p>1$ a série $p$ converge.
  Quando $p<1$ temos que $p-1<0$, então
  \[
   \lim_{b\to\infty} \frac{1}{\,b^{p-1}\,}
  \]
  diverge, portanto a integral diverge e consequentemente a série também.
\end{proof}
%------------------------------------------------------------------------------%

%------------------------------------------------------------------------------%
\begin{example}
  Verifique a convergência série
  \(\;
    \sum_{n=1}^\infty \frac{1}{n^2+1}
  \).

  \solution
  A série \(\sum_{n=1}^\infty \frac{1}{n^2+1}\) não é um série $p$,
  mas como $n^2+1>n^2$, segue que
  \[
      0 \;<\;\sum_{k=1}^n \frac{1}{k^2+1}
        \;<\; \sum_{k=1}^n \frac{1}{k^2}
      \qquad \forall\; n\in \N          \,.
  \]
  Entretanto, pela Proposição~\ref{pro:convergencia-serie-p}, temos que a série $p$
  \[
    \sum_{n=1}^\infty \frac{1}{n^2}
  \]
  converge, pois $p=2>1$. Dessa forma a sequência de somas parciais
  \[
    S_n=\sum_{k=1}^n \frac{1}{k^2+1}
  \]
  é limitada, e como os termos da série são não negativos,
  usamos o Corolário~\ref{cor:convergencia-serie-termos-nao-negativos}
  para verificar que a série converge.
\end{example}
%------------------------------------------------------------------------------%

%Veja outra solução para essa mesma série no Exemplo 4 da Seção 10.3 do
%livro do Thomas \cite{thomas_2}.
%A ideia é aplicar o teorema de forma mais direta, resolvendo a integral
%\[
%  \int_{1}^\infty \frac{1}{x^2+1}dx
%\]
%e concluindo que essa integral converge.

%------------------------------------------------------------------------------%
%------------------------------------------------------------------------------%
\subsection*{Estimativa de Erro}
%------------------------------------------------------------------------------%
%------------------------------------------------------------------------------%

Para analisarmos o erro cometido ao aproximarmos a soma de uma série $S$
por uma soma parcial $S_n$ vamos definir o erro ou resto da série.

%------------------------------------------------------------------------------%
\begin{definicao}{Resto da Série}{def:resto-serie}
  Dada uma série convergente, cuja soma é $S$
  \[
    S = \sum_{k=1}^\infty a_k
  \]
  podemos usar suas somas parciais
  \[
    S_n = \sum_{k=1}^n a_k
  \]
  como aproximações para $S$. Ao fazermos isso estamos cometendo
  um \emph{erro} igual ao \emph{Resto da Série}\index{Série!resto},
  isso é, o erro é igual a todos os termos que não foram somados
  \[
    R_n = S-S_n
        = \sum_{k=n+1}^\infty a_k       \,.
  \]
\end{definicao}
%------------------------------------------------------------------------------%

%------------------------------------------------------------------------------%
\begin{figure}
  \centering
  \begin{subfigure}{\linewidth}
    \centering
    \vspace{5mm}
    % Title: gl2ps_renderer figure
% Creator: GL2PS 1.4.2, (C) 1999-2020 C. Geuzaine
% For: Octave
% CreationDate: Sat Sep 18 16:46:57 2021
\setlength{\unitlength}{1pt}
\begin{picture}(0,0)
\includegraphics[width=16cm,height=5cm]{fig_teste_integral_erro_right}
\end{picture}%
\begin{picture}(453,141)(0,0)
\fontsize{10}{0}\selectfont\put(50.3937,10.9531){\makebox(0,0)[t]{\textcolor[rgb]{0.15,0.15,0.15}{{$n\!-\!1$}}}}
\fontsize{10}{0}\selectfont\put(85.6693,10.9531){\makebox(0,0)[t]{\textcolor[rgb]{0.15,0.15,0.15}{{$n\vphantom{1}$}}}}
\fontsize{10}{0}\selectfont\put(120.945,10.9531){\makebox(0,0)[t]{\textcolor[rgb]{0.15,0.15,0.15}{{$n\!+\!1$}}}}
\fontsize{10}{0}\selectfont\put(191.496,10.9531){\makebox(0,0)[t]{\textcolor[rgb]{0.15,0.15,0.15}{{$n\!+\!3$}}}}
\fontsize{10}{0}\selectfont\put(262.047,10.9531){\makebox(0,0)[t]{\textcolor[rgb]{0.15,0.15,0.15}{{$n\!+\!5$}}}}
\fontsize{10}{0}\selectfont\put(332.598,10.9531){\makebox(0,0)[t]{\textcolor[rgb]{0.15,0.15,0.15}{{$n\!+\!7$}}}}
\fontsize{10}{0}\selectfont\put(403.15,10.9531){\makebox(0,0)[t]{\textcolor[rgb]{0.15,0.15,0.15}{{$n\!+\!9$}}}}
\end{picture}

    \caption{Construindo os retângulos a direita a série é maior do que a integral.}
    \label{fig:teste_integral_erro_right}
  \end{subfigure}\\[2mm]
  \begin{subfigure}{\linewidth}
    \centering
    \input{fig_teste_integral_erro_left}
    \caption{Construindo os retângulos a esquerda a série é menor do que a integral.}
    \label{fig:teste_integral_erro_left}
  \end{subfigure}\\
  \caption{Estimativa do erro pelo testo da integral.}
  \label{fig:teste_integral_erro}
\end{figure}
%------------------------------------------------------------------------------%

Quando verificamos que uma série é convergente pelo teste da integral
podemos obter também estimativas para o erro cometido ao aproximarmos
a soma da série por uma das suas somas parciais.
Observe os gráficos da Figura~\ref{fig:teste_integral_erro} que
indicam o resto da série, $R_n$, após calcularmos a $n$-ésima soma parcial, $S_n$.
No Gráfico~\ref{fig:teste_integral_erro_right} os retângulos com área $a_k$
estão desenhados a direita do ponto $x=k$ e podemos perceber que
\[
  R_n > \int_{n+1}^{\infty} f(x)\,dx          \,.
\]
Fazendo a mesma comparação no Gráfico~\ref{fig:teste_integral_erro_left},
onde os retângulos estão a esquerda, vemos que
\[
  R_n < \int_{n}^{\infty} f(x)\,dx         \,.
\]
Note a mudança no início do intervalo de integração.
Temos assim limites inferior e superior para o erro cometido
ao aproximarmos a soma da série por uma soma parcial.
O próximo teorema formaliza esse resultado.

%------------------------------------------------------------------------------%
\begin{teorema}{Estimativa de Erro}{teo:estimativa-erro}
  Seja%
  \index{Série!estimativa de erro!teste integral}%
  \index{Estimativa de erro!teste integral}
  $a_k$ uma sequência tal que $a_k\geq 0$,
  para todo $k\geq N$, onde $N$ é um número natural.
  Seja $f$ uma função contínua, positiva e decrescente para $x\geq N$,
  satisfazendo $f(k)=a_k$, para todo $k\geq N$.
  Suponha que a série com termos $a_k$ converge para $S$, isso é,
  \[
    \sum_{k=1}^\infty a_k = S
  \]
  então o resto
  \[
    R_n=S-S_n
  \]
  satisfaz
  \[
    \int_{n+1}^\infty f(x)\,dx \;\leq\; R_n \;\leq\; \int_{n}^\infty f(x)\,dx        \,.
  \]
  Consequentemente,
  \[
    S_n+\int_{n+1}^\infty f(x)\,dx \;\leq\; S \;\leq\; S_n+\int_{n}^\infty f(x)\,dx        \,.
  \]
\end{teorema}
%------------------------------------------------------------------------------%
\begin{proof}
  A demonstração é parecida com a demonstração do Teorema~\ref{teo:teste-integral}.
  Basta reparar que $R_n$ é uma aproximação por baixo da integral
  \[
    \int_{n}^\infty f(x)\,dx
  \]
  já que representa a soma de retângulos de base 1 e altura
  $a_k=f(k)$, para $k=n+1, n+2, \dots$
  Por outro lado, $R_n$ é uma aproximação por cima da integral
  \[
    \int_{n+1}^\infty f(x)\,dx
  \]
  A última desigualdade segue somando $S_n$ na desigualdade anterior e usando que
  \[
    R_n = S - S_n              \,.
  \]
\end{proof}
%------------------------------------------------------------------------------%

%------------------------------------------------------------------------------%
\begin{example}
  Sabemos que a série
  \[
    \sum_{n=1}^\infty\frac{1}{\,n^2\,}
  \]
  converge, pois é uma série $p$, com $p=2$.
  Descubra $n$ tal que o resto $R_n$ da aproximação $S_n$ seja menor que \num{0.01}.

  \solution
  Sabemos que o Teste da Integral pode ser aplicado, pois
  \[
    \int_n^\infty\frac{\,dx\,}{x^2}
      = \lim_{b\to\infty}\left[ -\frac{1}{\,b\,} + \frac{1}{\,n\,} \right]
      = \frac{1}{\,n\,}                \,.
  \]
  Assim, pelo Teorema~\ref{teo:estimativa-erro},
  \[
    \frac{1}{\,n+1\,} \leq R_n \leq \frac{1}{\,n\,}           \,.
  \]
  Para garantir que $R_n<\num{0.01}$ impomos que
  \[
    \frac{1}{\,n\,} < \num{0.01}
  \]
  isolando $n$ optemos $n>100$. Concluímos então que a aproximação $S_{100}$
  tem um erro menor do que \num{0.01}.
\end{example}

%------------------------------------------------------------------------------%
